\phantomsection
\addcontentsline{toc}{section}{Résumé}
\begin{center}\bfseries Résumé\end{center}
\begingroup
\fontsize{9.5}{10.7}\selectfont
\setlength{\parskip}{1.5pt plus .5pt minus .5pt}
\noindent
On dérive les distributions de Bose--Einstein et de Fermi--Dirac, leur limite
diluée de Maxwell--Boltzmann et les lois de rayonnement de Planck,
Stefan--Boltzmann et Wien à partir des réalisations quantiques et thermiques
construites en holographie modulaire triadique et en mécanique dimensionnelle.
Le support arithmétique des chemins, appelé \emph{All Possible Paths}
(APP), est le tore discret \((\mathbb Z/9\mathbb Z)^2\), avec graphe de
Cayley additif de directions cardinales et deux évaluations, somme et
produit. La signature ternaire orientée, TRIT, prend ses valeurs dans
\(\{-1,0,+1\}\) et distingue l’orientation et le régime ; le noyau topologique
de phase, \emph{Topological Phase Kernel} (TPK), compose la sélection,
le transport, la mise à jour et le prolongement en conservant la mémoire.
Cette dynamique produit les états et les registres régionaux sur lesquels
l’évolution unitaire et l’action du calendrier sont définies.

La parité des réalisations détermine leurs complétions symétrique et
extérieure : la première admet des occupations entières non négatives ; la seconde
produit l’exclusion algébrique de Pauli. Les multiplicités déterminent les dénombrements microcanoniques ; la
factorisation de l’état de Gibbs donne les deux distributions d’équilibre,
dont la limite diluée est quantifiée au moyen de bornes explicites. Pour les
modes électromagnétiques transversaux libres, à potentiel chimique nul,
l’occupation bosonique et la densité de modes produisent le spectre de
Planck. On démontre la convergence des sommes discrètes au moyen de
bornes uniformes aux basses et hautes fréquences. Leurs moments déterminent
les densités de nombre, d’énergie et d’entropie ; l’intégration angulaire donne
le flux de Stefan--Boltzmann, et l’extrémisation des densités en
fréquence et en longueur d’onde donne les deux lois de Wien.

Le calendrier et l’holonomie produisent en outre un renouvellement de capacité
$(20,24,25,23)$, un porteur gradué de multiplicités
$1/380/649$ et un spectre géométrique natif de raison $1/10$. Son
raffinement tritique conserve le résidu et donne $1/1000$ ; la section marquée
$23/26/29$ évalue exactement
$1{,}380649\times10^{-23}$. L’origine, le marqueur et la provenance
énergétique sont transportés au moyen d’isométries explicites. Les courants de
bord fixent le déphasage sectoriel, et la composition avec l’action démontre
$G_a k_B\Theta_{{\rm clk},a}\log3=c^4L_\alpha$ et l’identité effective du
transducteur sur sa section positive.

L’entropie des configurations, la mesure des histoires et le spectre d’un
état réduit sont reliés par des applications de transport et de
réduction. Une identité de divergence relative caractérise la mesure
stationnaire d’entropie maximale de l’enveloppe surfacique--volumétrique.
La transduction thermique relie l’information, l’action et le calendrier à leurs
bases dimensionnelles. Les conditions de représentation, d’équilibre et de
convergence spécifient le domaine de chaque résultat.

Dans la réalisation finie \(\Lambda=(\mathbb Z/27\mathbb Z)^2\), la
transformée de Fourier unitaire établit
\(|\operatorname{supp}\psi|\allowbreak\,|\operatorname{supp}\widehat\psi|\ge729\)
et \(H_{\rm pos}(\psi)\allowbreak+H_F(\psi)\ge\log729\) ; quatre collections de sous-groupes
saturent les deux bornes. Dans le domaine dodécaphasique, les lecteurs de F044
distinguent des corrélations d’ordre que l’entropie marginale ne conserve pas.
L’observateur premier des vacances \(R_{\rm vac}\) compose l’incidence
sturmienne avec l’horloge logarithmique et satisfait une identité modale par
des sommes symétriques ou de Fejér. Ses six coupures jusqu’à \(10^8\) et les trente
modes enregistrés sont des témoins finis ; la reproduction indépendante
atteint \(10^6\). L’estimation asymptotique critique requiert des bornes modales
uniformes en \(k\).

\medskip
\noindent\textbf{Mots-clés :} Bose--Einstein ; Fermi--Dirac ; loi de Planck ;
états généalogiques ; complétion de Fock ; dynamique de Floquet ;
entropie d’occupation ; Hamiltonien modulaire ; limite thermodynamique ;
rayonnement thermique ; incertitude entropique ; vacances premières ; mémoire
ordonnée.
\par\endgroup
